\chapter{Rejillas: Redes Convolucionales}
\label{ch:grids}

\providecommand{\conv}{*}
\providecommand{\Ft}{\mathcal{F}}

La primera de las cinco Ges es tambi\'en la m\'as familiar. Una \emph{rejilla}
es el dominio euclidiano regular que subyace a las im\'agenes, las se\~nales de
audio y el v\'ideo: un conjunto de puntos $\Z^d$ sobre el que los datos se
muestrean en posiciones equiespaciadas. La simetr\'ia natural de un dominio as\'i
es la \emph{traslaci\'on}: desplazar la ret\'icula de muestreo por un vector
entero aplica la rejilla sobre s\'i misma. Las redes neuronales convolucionales
(\emph{convolutional neural networks}, CNN), las arquitecturas que encendieron la
revoluci\'on del aprendizaje profundo (\Cref{ch:introduction}), son precisamente
las capas que respetan esta simetr\'ia. Este cap\'itulo toma el esquema del
\Cref{ch:priors} ---convertir una simetr\'ia en una arquitectura--- y lo
instancia sobre la rejilla, recuperando la convoluci\'on como la \'unica
aplicaci\'on lineal equivariante a traslaciones y ensamblando luego toda la
maquinaria de \emph{pooling}, submuestreo y separaci\'on de escalas que hace
funcionar a las CNN profundas. Cerramos con una lectura espectral de la
convoluci\'on mediante el an\'alisis de Fourier, que a la vez explica qu\'e hacen
los filtros convolucionales y anticipa los m\'etodos espectrales sobre grafos y
variedades de cap\'itulos posteriores.

\section{Simetr\'ia de traslaci\'on}
\label{sec:grids:translation}

Las redes convolucionales tienen su origen en el reconocimiento cl\'asico de
patrones. La idea es de una sencillez desarmante: para detectar una
caracter\'istica en una imagen ---un borde horizontal, un borde vertical, una
esquina--- se desliza un peque\~no \emph{kernel} sobre la imagen y, en cada
posici\'on, se mide con qu\'e intensidad el parche local coincide con el kernel.
El avance del trabajo temprano de LeCun~\citep{lecun1989} consisti\'o en dejar de
dise\~nar estos kernels a mano y, en su lugar, \emph{aprenderlos} a partir de los
datos; apilar muchas de estas capas permite a la red construir desde trazos
simples hasta objetos abstractos, una jerarqu\'ia que culmin\'o en la victoria
decisiva de AlexNet en ImageNet~\citep{krizhevsky2012alexnet}.

\begin{example}[Los kernels de Sobel]
\label{ex:grids:sobel}
Un detector de bordes cl\'asico dise\~nado a mano es el par de \emph{kernels de
Sobel}, aproximaciones discretas de las derivadas parciales de la intensidad de
la imagen $I(x,y)$:
\[
K_x = \begin{pmatrix} -1 & 0 & 1 \\ -2 & 0 & 2 \\ -1 & 0 & 1 \end{pmatrix},
\qquad
K_y = \begin{pmatrix} -1 & -2 & -1 \\ 0 & 0 & 0 \\ 1 & 2 & 1 \end{pmatrix},
\]
que aproximan $\partial I/\partial x$ y $\partial I/\partial y$, con intensidad
de borde $\norm{\nabla I} \approx \sqrt{(I \conv K_x)^2 + (I \conv K_y)^2}$. Los
kernels revelan un tema recurrente: los detectores de bordes son
\emph{antisim\'etricos} a lo largo de la direcci\'on de detecci\'on y suman cero
(eliminando la componente constante), mientras que los kernels de suavizado son
\emph{sim\'etricos}. Estas simetr\'ias algebraicas reflejan las transformaciones
geom\'etricas a las que responde cada kernel.
\end{example}

La propiedad que define a una capa convolucional aprendida es la
\emph{equivarianza a traslaciones}: si la entrada se desplaza, la salida se
desplaza exactamente de la misma manera. Una caracter\'istica debe reconocerse
con independencia de d\'onde aparezca en la imagen, y su respuesta debe moverse
con ella. Formalmente, el grupo que act\'ua sobre los datos en rejilla es el
grupo de traslaciones enteras $\Z^d$, que act\'ua sobre una se\~nal
$f \colon \Z^d \to \R$ mediante
\[
  (T_b f)(x) = f(x - b), \qquad b \in \Z^d .
\]
Esta es la instancia no trivial m\'as simple del principio de equivarianza del
\Cref{ch:priors}: entre todas las aplicaciones lineales concebibles de se\~nales
en se\~nales, exigir la conmutaci\'on con todo $T_b$ reduce el espacio de
dise\~no a una \'unica familia ---las convoluciones---.

\section{La convoluci\'on como aplicaci\'on lineal equivariante}
\label{sec:grids:convolution}

Comenzamos por la operaci\'on en s\'i. En la pr\'actica del aprendizaje profundo,
la capa aplicada a una imagen es
\begin{equation}
  (I \conv K)(x,y) = \sum_{u=-k}^{k}\sum_{v=-k}^{k} I(x+u,\, y+v)\, K(u,v),
  \label{eq:grids:crosscorr}
\end{equation}
que, en rigor, es una \emph{correlaci\'on cruzada} y no una convoluci\'on: la
convoluci\'on matem\'atica lleva un cambio de signo, $I(x-u,\,y-v)$. Como los
pesos del kernel se aprenden, la distinci\'on es irrelevante ---la red
simplemente aprende el kernel reflejado--- y seguimos la convenci\'on universal
de llamar convoluci\'on a la \Cref{eq:grids:crosscorr}. La misma receta se
extiende a cualquier dimensi\'on: una versi\'on unidimensional procesa audio y
series temporales, y una tridimensional procesa vol\'umenes y v\'ideo.

\begin{figure}[t]
\centering
\begin{tikzpicture}[
    cell/.style={minimum size=0.62cm, inner sep=0pt, outer sep=0pt,
                 draw=gdlgray!50, line width=0.3pt},
    inputcell/.style={cell, fill=gdlblue!15},
    kernelhighlight/.style={cell, fill=gdlorange!25, draw=gdlorange!60, line width=0.5pt},
    kernelcell/.style={cell, fill=gdlorange!20, draw=gdlorange!60, line width=0.5pt},
    outputcell/.style={cell, fill=gdlgreen!15, draw=gdlgray!50},
    outputhighlight/.style={cell, fill=gdlgreen!30, draw=gdlgreen!60, line width=0.6pt},
    lbl/.style={font=\small\bfseries},
]
\def\cs{0.62}
\def\ix{0}\def\iy{0}
\foreach \r [count=\ri from 0] in {{1,0,2,1,3},{2,1,0,3,1},{0,3,1,2,0},{1,2,3,0,1},{3,0,1,2,2}} {
  \foreach \val [count=\ci from 0] in \r {
    \pgfmathsetmacro{\px}{\ix + \ci*\cs + \cs/2}
    \pgfmathsetmacro{\py}{\iy + (4-\ri)*\cs + \cs/2}
    \node[inputcell] at (\px,\py) {\scriptsize\color{gdlblue!70!black}\val};
  }}
\foreach \r [count=\ri from 0] in {{1,0,2},{2,1,0},{0,3,1}} {
  \foreach \val [count=\ci from 0] in \r {
    \pgfmathsetmacro{\px}{\ix + \ci*\cs + \cs/2}
    \pgfmathsetmacro{\py}{\iy + (4-\ri)*\cs + \cs/2}
    \node[kernelhighlight] at (\px,\py) {\scriptsize\color{gdlblue!70!black}\val};
  }}
\draw[gdlorange!80, line width=1.3pt, rounded corners=1pt] (\ix,\iy+2*\cs) rectangle ++(3*\cs,3*\cs);
\node[lbl, text=gdlblue!70!black] at (\ix+2.5*\cs,\iy-0.5*\cs) {Entrada $\mathbf{X}$};
\pgfmathsetmacro{\starx}{\ix + 5*\cs + 0.5}
\pgfmathsetmacro{\stary}{\iy + 2.5*\cs}
\node[font=\large, text=gdlgray!80!black] at (\starx,\stary) {$\star$};
\pgfmathsetmacro{\kx}{\starx + 0.5}
\pgfmathsetmacro{\ky}{\stary - 1.5*\cs}
\foreach \r [count=\ri from 0] in {{1,0,-1},{1,0,-1},{1,0,-1}} {
  \foreach \val [count=\ci from 0] in \r {
    \pgfmathsetmacro{\px}{\kx + \ci*\cs + \cs/2}
    \pgfmathsetmacro{\py}{\ky + (2-\ri)*\cs + \cs/2}
    \node[kernelcell] at (\px,\py) {\scriptsize\color{gdlorange!70!black}\val};
  }}
\draw[gdlorange!80, line width=1.1pt, rounded corners=1pt] (\kx,\ky) rectangle ++(3*\cs,3*\cs);
\node[lbl, text=gdlorange!70!black] at (\kx+1.5*\cs,\ky-0.5*\cs) {Kernel $\mathbf{W}$};
\pgfmathsetmacro{\eqx}{\kx + 3*\cs + 0.5}
\node[font=\large, text=gdlgray!80!black] at (\eqx,\stary) {$=$};
\pgfmathsetmacro{\ox}{\eqx + 0.5}
\pgfmathsetmacro{\oy}{\ky}
\foreach \r [count=\ri from 0] in {{0,-2,-1},{-1,1,2},{-1,1,2}} {
  \foreach \val [count=\ci from 0] in \r {
    \pgfmathsetmacro{\px}{\ox + \ci*\cs + \cs/2}
    \pgfmathsetmacro{\py}{\oy + (2-\ri)*\cs + \cs/2}
    \ifnum\ri=0 \ifnum\ci=0
      \node[outputhighlight] at (\px,\py) {\scriptsize\bfseries\color{gdlgreen!70!black}\val};
    \else
      \node[outputcell] at (\px,\py) {\scriptsize\color{gdlgreen!70!black}\val};
    \fi \else
      \node[outputcell] at (\px,\py) {\scriptsize\color{gdlgreen!70!black}\val};
    \fi
  }}
\draw[gdlgreen!60, line width=0.8pt, rounded corners=1pt] (\ox,\oy) rectangle ++(3*\cs,3*\cs);
\node[lbl, text=gdlgreen!70!black] at (\ox+1.5*\cs,\oy-0.5*\cs) {Salida $(\mathbf{X} \star \mathbf{W})$};
\pgfmathsetmacro{\arcsrcx}{\ix + 1.5*\cs}
\pgfmathsetmacro{\arcsrcy}{\iy + 5*\cs + 0.08}
\pgfmathsetmacro{\arctgtx}{\ox + 0.5*\cs}
\pgfmathsetmacro{\arctgty}{\oy + 3*\cs + 0.08}
\draw[gdlpurple!55, dashed, line width=0.7pt, -stealth, shorten >=1pt]
  (\arcsrcx,\arcsrcy) to[out=40,in=140] (\arctgtx,\arctgty);
\end{tikzpicture}
\caption[La operaci\'on de convoluci\'on 2D]{La operaci\'on de convoluci\'on
bidimensional. Un kernel $\mathbf{W}$ de $3 \times 3$ se desliza sobre la entrada
$\mathbf{X}$, y en cada posici\'on el producto interno local
$\sum_{u,v} X(x+u,y+v)\,W(u,v)$ produce una entrada del mapa de activaci\'on de
salida $(\mathbf{X} \star \mathbf{W})$. La celda de salida resaltada vale
$1{\cdot}1 + 0{\cdot}0 + 2{\cdot}({-}1) + 2{\cdot}1 + 1{\cdot}0 + 0{\cdot}({-}1)
+ 0{\cdot}1 + 3{\cdot}0 + 1{\cdot}({-}1) = 0$.}
\label{fig:grids:conv}
\end{figure}

Para la teor\'ia resulta m\'as limpio trabajar con la convoluci\'on genuina sobre
el dominio completo, que definimos a continuaci\'on.

\begin{definition}[Convoluci\'on discreta]
\label{def:grids:discrete-conv}
Para funciones $f, g \colon \Z^d \to \R$, la \emph{convoluci\'on discreta} es
\[
  (f \conv g)(x) = \sum_{a \in \Z^d} f(a)\, g(x - a) .
\]
Cuando $f$ tiene soporte finito, como en una imagen digital, la suma se reduce a
los finitos puntos en que $f(a) \neq 0$.
\end{definition}

\begin{definition}[Convoluci\'on continua]
\label{def:grids:continuous-conv}
Para funciones $f, g \colon \R^d \to \R$, la \emph{convoluci\'on continua} es
\[
  (f \conv g)(x) = \int_{\R^d} f(a)\, g(x - a)\, \mathrm{d}a,
\]
bien definida cuando $f, g \in L^1(\R^d)$ o bajo hip\'otesis de soporte
compacto.
\end{definition}

La versi\'on discreta es una discretizaci\'on por sumas de Riemann de la
continua; en el procesamiento de imagen digital usamos la forma discreta,
mientras que las se\~nales continuas y el an\'alisis te\'orico emplean la
integral. La convoluci\'on posee una rica estructura algebraica que sustenta todo
lo que sigue.

\begin{proposition}[Propiedades algebraicas de la convoluci\'on]
\label{prop:grids:algebra}
Para $f, g, h \colon \Z^d \to \R$ de soporte finito:
\begin{enumerate}
  \item \textbf{Conmutatividad:} $f \conv g = g \conv f$.
  \item \textbf{Asociatividad:} $(f \conv g) \conv h = f \conv (g \conv h)$.
  \item \textbf{Distributividad:} $f \conv (g + h) = f \conv g + f \conv h$.
  \item \textbf{Elemento neutro:} $f \conv \delta = f$, donde $\delta$ es la
    delta de Kronecker.
\end{enumerate}
\end{proposition}

\begin{proof}
La conmutatividad se sigue del cambio de variable $b = x - a$, que transforma
$\sum_a f(a)g(x-a)$ en $\sum_b f(x-b)g(b)$. La asociatividad se obtiene
desarrollando ambos miembros, intercambiando el orden de sumaci\'on (un Fubini
discreto) y sustituyendo $c = b - a$. La distributividad es inmediata por la
linealidad de la suma. El elemento neutro se cumple porque $\delta(x-a)$ es no
nulo solo en $a = x$, de modo que $\sum_a f(a)\delta(x-a) = f(x)$.
\end{proof}

Estas propiedades tienen consecuencias arquitect\'onicas directas. La
asociatividad implica que apilar dos capas convolucionales (ignorando las no
linealidades) es a su vez una convoluci\'on, de modo que la profundidad compone
limpiamente; la conmutatividad y el elemento neutro hacen que el \'algebra de
filtros se comporte tan bien como la multiplicaci\'on ordinaria. El resultado
central, sin embargo, es el v\'inculo con la simetr\'ia.

\begin{theorem}[Equivarianza a traslaciones de la convoluci\'on]
\label{thm:grids:equivariance}
La convoluci\'on es equivariante al grupo de traslaciones $\Z^d$: para todo
$b \in \Z^d$,
\[
  T_b(f \conv g) = (T_b f) \conv g = f \conv (T_b g),
\]
donde $(T_b f)(x) = f(x-b)$.
\end{theorem}

\begin{proof}
Para la primera igualdad, sustituimos $a' = a - b$:
\[
  \big[(T_b f) \conv g\big](x) = \sum_a f(a-b)\,g(x-a)
   = \sum_{a'} f(a')\,g\big((x-b)-a'\big) = (f \conv g)(x-b).
\]
Para la segunda, $(f\conv g)(x-b) = \sum_a f(a)\,g\big((x-a)-b\big)
= \sum_a f(a)\,(T_b g)(x-a) = \big[f \conv (T_b g)\big](x)$.
\end{proof}

Una peque\~na instancia num\'erica hace tangible el teorema.

\begin{example}[Equivarianza en una dimensi\'on]
\label{ex:grids:equiv-1d}
Tomemos el kernel detector de bordes $K = [-1, 0, 1]$ y la se\~nal
$f = [1, 2, 1, 0, 0]$, con relleno de ceros fuera del dominio. La convoluci\'on
da $f \conv K = [2, 0, -2, -1, 0]$. Si primero desplazamos,
$T_1 f = [0, 1, 2, 1, 0]$, y luego convolucionamos, obtenemos
$(T_1 f) \conv K = [1, 2, 0, -2, -1]$, que es exactamente $T_1(f \conv K)$.
Desplazar y filtrar, o filtrar y desplazar: el resultado es el mismo.
\end{example}

El rec\'iproco es lo que convierte a la convoluci\'on en \emph{la} capa de la
rejilla, y no simplemente en \emph{una} capa: cualquier aplicaci\'on lineal entre
se\~nales sobre $\Z^d$ que conmute con todas las traslaciones es necesariamente
una convoluci\'on. El razonamiento es breve. Un operador lineal equivariante a
traslaciones $\Phi$ queda determinado por su respuesta al impulso unitario
$\delta$; ll\'amese a esa respuesta $w = \Phi\delta$. Toda se\~nal se descompone
como $f = \sum_a f(a)\, T_a \delta$ y, por linealidad y equivarianza,
$\Phi f = \sum_a f(a)\, T_a w = f \conv w$. As\'i, la equivarianza a la
traslaci\'on no solo \emph{permite} la convoluci\'on ---la \emph{obliga}---. Este
es el esquema del \Cref{ch:priors} realizado en su forma m\'as pura, y explica la
se\~na de identidad arquitect\'onica de las CNN.

\paragraph{Compartici\'on de pesos y localidad.}
El \Cref{thm:grids:equivariance} admite una lectura concreta: el mismo kernel
$w$ se aplica en cada posici\'on espacial. Esta \emph{compartici\'on de pesos}
(\emph{weight sharing}) es la cara algebraica de la equivarianza. Frente a una
capa totalmente conectada, que aprender\'ia un peso independiente para cada par
de posiciones de entrada y salida, una convoluci\'on comparte un \'unico kernel
peque\~no en toda la imagen, reduciendo el n\'umero de par\'ametros de
cuadr\'atico en el n\'umero de p\'ixeles al tama\~no constante del kernel.
Restringir el kernel a una ventana peque\~na a\~nade \emph{localidad}: cada
salida depende solo de un vecindario local. Apilar tales capas hace crecer
progresivamente el \emph{campo receptivo}, de modo que las primeras capas
capturan trazos simples y las m\'as profundas los componen en estructuras cada
vez m\'as abstractas.

\paragraph{Los l\'imites de la traslaci\'on.}
La fortaleza del sesgo inductivo de las CNN es tambi\'en su frontera. Una red
convolucional est\'andar es equivariante a traslaciones y a nada m\'as. Si se
rota el d\'igito ``6'' $180^\circ$, se convierte en un ``9''; una CNN no dispone
de un mecanismo intr\'inseco para relacionar ambos y debe aprender las variantes
rotadas mediante aumentaci\'on de datos. En general $f(R\,x) \neq R\, f(x)$ para
una rotaci\'on $R$. Eliminar esta restricci\'on ---ampliar el grupo de simetr\'ia
m\'as all\'a de la traslaci\'on para incluir rotaciones y reflexiones y,
finalmente, abandonar por completo la rejilla global--- es el hilo que recorre
las Ges restantes: redes equivariantes a grupos (\Cref{ch:groups}), grafos
(\Cref{ch:graphs}), variedades (\Cref{ch:geodesics}) y gauges
(\Cref{ch:gauges}).

\section{\emph{Pooling}, submuestreo y separaci\'on de escalas}
\label{sec:grids:pooling}

La convoluci\'on por s\'i sola no basta para construir un sistema de
reconocimiento profundo. Hace falta un segundo ingrediente: una forma de
engrosar la rejilla a medida que la informaci\'on asciende por la red, de modo
que las caracter\'isticas locales detectadas a resoluci\'on fina se agreguen en
caracter\'isticas globales a resoluci\'on gruesa. Este es el papel del
\emph{pooling} y del \emph{submuestreo}.

Una operaci\'on de \emph{pooling} divide la rejilla en peque\~nos bloques
---t\'ipicamente $2 \times 2$--- y sustituye cada bloque por un \'unico valor
resumen: el m\'aximo (\emph{max pooling}) o el promedio (\emph{average pooling})
de sus entradas. La salida vive en una rejilla de la mitad de lado, de modo que
la resoluci\'on espacial cae en un factor cuatro mientras que el campo receptivo
de cada kernel posterior, medido en p\'ixeles originales, se duplica. La
convoluci\'on con paso (\emph{strided convolution}) logra el mismo engrosamiento
dentro de la propia convoluci\'on, evaluando la salida solo cada $s$ posiciones.

El \emph{pooling} es \emph{invariante a traslaciones a la escala del bloque}:
desplazar la entrada menos que el tama\~no del bloque deja la salida (de forma
aproximada) inalterada. Esta es exactamente la robustez local que se desea de un
detector de caracter\'isticas ---el p\'ixel exacto en que se sit\'ua un borde no
deber\'ia importar, solo que haya un borde en las cercan\'ias--- y complementa la
equivarianza global de la convoluci\'on. Apilar convoluci\'on, no linealidad y
\emph{pooling} produce el bloque can\'onico de las CNN, y repetirlo construye la
pir\'amide multiescala en la que las primeras capas ven detalle local fino y las
posteriores ven estructura global gruesa.

\paragraph{Separaci\'on de escalas.}
La justificaci\'on m\'as profunda de esta arquitectura es el sesgo geom\'etrico
de la \emph{separaci\'on de escalas}. Las se\~nales naturales son
aproximadamente composicionales: un patr\'on complejo se construye a partir de
patrones m\'as simples a escalas m\'as finas, y las interacciones entre escalas
muy distintas son d\'ebiles. Las caras se componen de ojos y bocas, que se
componen de bordes y texturas, que se componen de cambios locales de intensidad.
Una funci\'on con esta estructura puede aproximarse componiendo operadores
\emph{locales} (convoluciones) con operadores de \emph{engrosamiento}
(\emph{pooling}), evitando el coste exponencial de modelar todas las
interacciones de largo alcance a la vez. La simetr\'ia de traslaci\'on y la
separaci\'on de escalas juntas ---compartici\'on de pesos entre posiciones,
engrosamiento jer\'arquico entre escalas--- constituyen los dos sesgos
geom\'etricos que hacen a las redes convolucionales a la vez estad\'isticamente
eficientes y computacionalmente tratables. Son la especializaci\'on al dominio de
la rejilla del esquema general, y reaparecen, debidamente generalizados, en cada
cap\'itulo posterior.

\section{Una visi\'on espectral: an\'alisis arm\'onico de la convoluci\'on}
\label{sec:grids:spectral}

Existe una segunda manera, dual, de entender la convoluci\'on, y se convertir\'a
en el principio organizador de la convoluci\'on sobre grafos y variedades. La
idea es pasar del dominio espacial al dominio \emph{frecuencial} mediante la
transformada de Fourier.

\begin{definition}[Transformada de Fourier]
\label{def:grids:fourier}
Para $f \colon \R^d \to \C$ integrable, la \emph{transformada de Fourier} es
\[
  \Ft[f](\omega) = \hat{f}(\omega) = \int_{\R^d} f(x)\, e^{-2\pi i\, \omega \cdot x}\, \mathrm{d}x,
  \qquad \omega \in \R^d,
\]
con inversa $f(x) = \int_{\R^d} \hat{f}(\omega)\, e^{2\pi i\, \omega \cdot x}\, \mathrm{d}\omega$,
v\'alida (por ejemplo) cuando $f, \hat{f} \in L^1(\R^d)$.
\end{definition}

Las exponenciales complejas $e^{2\pi i\,\omega\cdot x}$ son exactamente las
\emph{autofunciones} de la traslaci\'on: desplazar el argumento multiplica la
exponencial por una fase. Esto es lo que hace del an\'alisis de Fourier el
lenguaje natural de la simetr\'ia de traslaci\'on ---y anticipa el papel central
de la base de autovectores del laplaciano sobre grafos (\Cref{ch:graphs}), que
desempe\~na la misma funci\'on en dominios irregulares---. La recompensa es el
teorema de convoluci\'on.

\begin{theorem}[Teorema de convoluci\'on]
\label{thm:grids:conv-theorem}
Para $f, g \colon \R^d \to \C$ adecuadas,
\[
  \Ft[f \conv g] = \Ft[f] \cdot \Ft[g] .
\]
La convoluci\'on en el dominio espacial es multiplicaci\'on puntual en el dominio
frecuencial.
\end{theorem}

El teorema convierte la operaci\'on global y deslizante de la convoluci\'on en un
reescalado independiente de cada componente frecuencial. Tres consecuencias
importan aqu\'i. Primera, es la base de la convoluci\'on r\'apida: mediante la
transformada r\'apida de Fourier (\emph{Fast Fourier Transform}, FFT), una
convoluci\'on que cuesta $O(N^2)$ de forma directa puede calcularse en
$O(N \log N)$ transformando, multiplicando y antitransformando. Segunda, ofrece
una interpretaci\'on limpia de lo que \emph{hace} un filtro: un kernel act\'ua
como una \emph{m\'ascara frecuencial}, amplificando unas frecuencias y atenuando
otras. Tercera, revela el significado de los kernels dise\~nados a mano del
\Cref{ex:grids:sobel}. Un kernel de suavizado (gaussiano) tiene una transformada
de Fourier concentrada cerca de $\omega = 0$, por lo que es un filtro
\emph{paso bajo} que suprime el detalle fino; un detector de bordes de Sobel o
laplaciano tiene una transformada que se anula en $\omega = 0$ y crece con
$\norm{\omega}$, por lo que es un filtro \emph{paso alto} que responde a los
cambios r\'apidos de intensidad. Aprender un kernel es, pues, aprender una
respuesta en frecuencia.

\begin{example}[Localizaci\'on dual]
\label{ex:grids:gaussian-pair}
La gaussiana es su propia transformada de Fourier salvo reescalado:
$\Ft[e^{-\alpha x^2}](\omega) = \sqrt{\pi/\alpha}\; e^{-\pi^2 \omega^2/\alpha}$.
Una gaussiana \emph{estrecha} en el espacio (poca extensi\'on espacial,
$\alpha$ grande) tiene una transformada \emph{ancha} en frecuencia, y viceversa.
Este compromiso espacio--frecuencia es el principio de incertidumbre, y tiene un
significado arquitect\'onico directo: un filtro no puede estar simult\'aneamente
muy localizado en el espacio y ser muy selectivo en frecuencia. Las CNN resuelven
la tensi\'on a lo largo de la profundidad, usando kernels locales peque\~nos a
escalas finas y, mediante el \emph{pooling}, kernels efectivamente mayores y m\'as
selectivos en frecuencia a escalas gruesas.
\end{example}

La perspectiva espectral aclara tambi\'en la relaci\'on entre las dos mitades de
este cap\'itulo. La convoluci\'on es el operador equivariante a traslaciones en
el dominio espacial; en el dominio espectral es el operador \emph{diagonalizado}
por la base de Fourier. Ambas descripciones dicen lo mismo sobre la rejilla, pero
la espectral sobrevive a la p\'erdida de la rejilla. En un grafo no hay noci\'on
de ``deslizar un kernel'', pero s\'i hay un laplaciano cuyos autovectores
generalizan la base de Fourier, y puede definirse la convoluci\'on como
multiplicaci\'on en esa base. Este es el puente de las rejillas a los grafos que
el \Cref{ch:graphs} desarrolla por completo, y muestra que la humilde capa
convolucional es la primera instancia de una construcci\'on que se extiende a
trav\'es de las cinco Ges.

\section{Caso de estudio: CNN para im\'agenes y audio}
\label{sec:grids:case-study}

Conviene ver c\'omo se ensamblan las piezas en las arquitecturas que definieron
el campo. AlexNet~\citep{krizhevsky2012alexnet}, la red que gan\'o la
competici\'on ImageNet de 2012 y redujo su error top-5 del $26{,}2\%$ al
$15{,}3\%$, es una pila de ocho capas aprendidas: cinco capas convolucionales
intercaladas con \emph{max pooling}, seguidas de tres capas totalmente conectadas
y un clasificador final. Las primeras capas convolucionales, que operan sobre la
imagen a resoluci\'on completa, aprenden detectores de bordes orientados y de
color que recuerdan llamativamente a los kernels de Sobel dise\~nados a mano del
\Cref{ex:grids:sobel}; las sucesivas etapas de \emph{pooling} engrosan la rejilla
de modo que las capas m\'as profundas, con campos receptivos efectivos mayores,
responden a texturas, partes y, finalmente, objetos completos. La equivarianza a
traslaciones de las convoluciones garantiza que un objeto se reconozca all\'i
donde aparezca, y el \emph{pooling} aporta la robustez local frente a peque\~nos
desplazamientos. El \'exito se apoy\'o en los tres factores destacados en el
\Cref{ch:introduction}: grandes conjuntos de datos etiquetados, c\'omputo en GPU
y una arquitectura ---la pila convolucional--- cuyo sesgo inductivo encaja con la
geometr\'ia de las im\'agenes.

La misma construcci\'on se transfiere al audio con un cambio de dimensi\'on. Una
forma de onda o su espectrograma es una rejilla uni- o bidimensional con
simetr\'ia de traslaci\'on a lo largo del eje temporal: un evento sonoro debe
detectarse con independencia de cu\'ando ocurra. Las convoluciones
unidimensionales sobre la forma de onda cruda, o las bidimensionales sobre el
espectrograma tiempo--frecuencia, comparten pesos a lo largo del tiempo
exactamente igual que las convoluciones de imagen comparten pesos en el espacio,
con el \emph{pooling} agregando caracter\'isticas ac\'usticas locales en
estructura de mayor alcance. La visi\'on espectral de la
\Cref{sec:grids:spectral} es especialmente natural aqu\'i, pues el espectrograma
es ya una representaci\'on frecuencial y los filtros convolucionales act\'uan
sobre \'el como detectores paso banda aprendidos.

Estos ejemplos hacen concreta la tesis del cap\'itulo. Las redes convolucionales
no son un truco de ingenier\'ia ad hoc, sino la realizaci\'on can\'onica de un
\'unico principio ---respetar la simetr\'ia de traslaci\'on de la rejilla---
combinado con el sesgo de separaci\'on de escalas que legitima el engrosamiento
jer\'arquico. Mantener fijo el principio y ampliar el grupo de simetr\'ia es
precisamente el programa de los cap\'itulos que siguen, comenzando en el
\Cref{ch:groups} con redes equivariantes a rotaciones y reflexiones.
